% ========================================================
% BÀI TẬP TOÁN 5 - KẾT NỐI TRI THỨC VỚI CUỘC SỐNG
% Bài 2: Ôn tập các phép tính với số tự nhiên
% Định dạng: BTTeX Part III - Trắc nghiệm điền số ngắn
% ========================================================

% Câu 1
\begin{ex}
  Tính kết quả của phép cộng: $48\,725 + 23\,618$.
  \shortans{72343}
  \loigiai{
    Thực hiện đặt tính và tính: $48\,725 + 23\,618 = 72\,343$.
  }
\end{ex}

% Câu 2
\begin{ex}
  Tính kết quả của phép trừ: $85\,400 - 37\,850$.
  \shortans{47550}
  \loigiai{
    Thực hiện đặt tính trừ: $85\,400 - 37\,850 = 47\,550$.
  }
\end{ex}

% Câu 3
\begin{ex}
  Tính giá trị của phép nhân: $308 \times 25$.
  \shortans{7700}
  \loigiai{
    Thực hiện phép nhân: $308 \times 25 = 7\,700$.
  }
\end{ex}

% Câu 4
\begin{ex}
  Tính thương của phép chia: $9\,720 : 36$.
  \shortans{270}
  \loigiai{
    Thực hiện phép chia: $9\,720 : 36 = 270$.
  }
\end{ex}

% Câu 5
\begin{ex}
  Tìm số dư của phép chia: $15\,872 : 45$.
  \shortans{32}
  \loigiai{
    Thực hiện phép chia $15\,872 : 45 = 352$ (dư $32$). Vậy số dư là 32.
  }
\end{ex}

% Câu 6
\begin{ex}
  Tính nhanh giá trị biểu thức bằng cách hợp lí: $25 \times 137 \times 4$.
  \shortans{13700}
  \loigiai{
    Áp dụng tính chất giao hoán và kết hợp: $(25 \times 4) \times 137 = 100 \times 137 = 13\,700$.
  }
\end{ex}

% Câu 7
\begin{ex}
  Tính nhanh giá trị biểu thức: $48 \times 65 + 48 \times 35$.
  \shortans{4800}
  \loigiai{
    Áp dụng tính chất phân phối: $48 \times (65 + 35) = 48 \times 100 = 4\,800$.
  }
\end{ex}

% Câu 8
\begin{ex}
  Tính nhanh giá trị biểu thức: $125 \times 79 - 125 \times 71$.
  \shortans{1000}
  \loigiai{
    Ta có: $125 \times (79 - 71) = 125 \times 8 = 1\,000$.
  }
\end{ex}

% Câu 9
\begin{ex}
  Tìm số tự nhiên $x$ biết: $x \times 28 = 4\,200$.
  \shortans{150}
  \loigiai{
    $x = 4\,200 : 28 = 150$.
  }
\end{ex}

% Câu 10
\begin{ex}
  Tìm số tự nhiên $x$ biết: $9\,600 : x = 32$.
  \shortans{300}
  \loigiai{
    $x = 9\,600 : 32 = 300$.
  }
\end{ex}

% Câu 11
\begin{ex}
  Tìm số tự nhiên $x$ biết: $x : 15 + 85 = 145$.
  \shortans{900}
  \loigiai{
    Ta có $x : 15 = 145 - 85 = 60 \implies x = 60 \times 15 = 900$.
  }
\end{ex}

% Câu 12
\begin{ex}
  Tính giá trị biểu thức: $2\,500 - (180 + 320) \times 3$.
  \shortans{1000}
  \loigiai{
    Thực hiện trong ngoặc trước: $180 + 320 = 500$. Nhân với 3 được $1\,500$. Kết quả: $2\,500 - 1\,500 = 1\,000$.
  }
\end{ex}

% Câu 13
\begin{ex}
  Hai kho thóc chứa tất cả $450$ tấn thóc. Kho thứ nhất chứa nhiều hơn kho thứ hai $50$ tấn thóc. Hỏi kho thứ nhất chứa bao nhiêu tấn thóc?
  \shortans{250}
  \loigiai{
    Kho thứ nhất chứa số tấn thóc là: $(450 + 50) : 2 = 250$ (tấn).
  }
\end{ex}

% Câu 14
\begin{ex}
  Bốn tổ công nhân sản xuất được số sản phẩm lần lượt là $120, 135, 145$ và $160$ sản phẩm. Hỏi trung bình mỗi tổ sản xuất được bao nhiêu sản phẩm?
  \shortans{140}
  \loigiai{
    Trung bình mỗi tổ sản xuất được: $(120 + 135 + 145 + 160) : 4 = 560 : 4 = 140$ (sản phẩm).
  }
\end{ex}

% Câu 15
\begin{ex}
  Mua $6$ quyển vở cùng loại hết $54\,000$ đồng. Hỏi mua $9$ quyển vở như thế hết bao nhiêu nghìn đồng?
  \shortans{81}
  \loigiai{
    Giá 1 quyển vở là $54\,000 : 6 = 9\,000$ đồng. Mua 9 quyển vở hết $9\,000 \times 9 = 81\,000$ đồng, tức 81 nghìn đồng.
  }
\end{ex}

% Câu 16
\begin{ex}
  Một ô tô đi quãng đường $180\text{ km}$ hết $4$ giờ. Hỏi với vận tốc đó, trong $6$ giờ ô tô đi được bao nhiêu ki-lô-mét?
  \shortans{270}
  \loigiai{
    Mỗi giờ ô tô đi được: $180 : 4 = 45\text{ km}$. Trong 6 giờ ô tô đi được: $45 \times 6 = 270\text{ km}$.
  }
\end{ex}

% Câu 17
\begin{ex}
  Tính giá trị biểu thức: $120 \times 5 - 450 : 9 + 35$.
  \shortans{585}
  \loigiai{
    $120 \times 5 = 600$; $450 : 9 = 50$. Biểu thức bằng $600 - 50 + 35 = 585$.
  }
\end{ex}

% Câu 18
\begin{ex}
  Tìm số bị chia của phép chia có số chia là $35$, thương là $24$ và số dư là số dư lớn nhất có thể.
  \shortans{874}
  \loigiai{
    Số dư lớn nhất có thể là $35 - 1 = 34$. Số bị chia là: $35 \times 24 + 34 = 840 + 34 = 874$.
  }
\end{ex}

% Câu 19
\begin{ex}
  Một hình chữ nhật có chu vi $140\text{ m}$, chiều dài hơn chiều rộng $10\text{ m}$. Diện tích hình chữ nhật đó là bao nhiêu mét vuông?
  \shortans{1200}
  \loigiai{
    Nửa chu vi là $140 : 2 = 70\text{ m}$. Chiều dài là $(70 + 10) : 2 = 40\text{ m}$. Chiều rộng là $70 - 40 = 30\text{ m}$. Diện tích là $40 \times 30 = 1\,200\text{ m}^2$.
  }
\end{ex}

% Câu 20
\begin{ex}
  Tìm số tự nhiên lớn nhất có ba chữ số chia hết cho cả $2, 5$ và $9$.
  \shortans{990}
  \loigiai{
    Số chia hết cho cả 2 và 5 phải có tận cùng là 0. Để số chia hết cho 9 thì tổng các chữ số chia hết cho 9. Số lớn nhất có ba chữ số thỏa mãn là $990$.
  }
\end{ex}

% Câu 21
\begin{ex}
  Hiện nay mẹ hơn con $26$ tuổi. Tổng số tuổi của hai mẹ con là $46$ tuổi. Hỏi sau $2$ năm nữa tuổi của con là bao nhiêu?
  \shortans{12}
  \loigiai{
    Tuổi con hiện nay là: $(46 - 26) : 2 = 10$ tuổi. Sau 2 năm nữa tuổi con là $10 + 2 = 12$ tuổi.
  }
\end{ex}

% Câu 22
\begin{ex}
  Tính giá trị biểu thức: $1 + 2 + 3 + \dots + 20$.
  \shortans{210}
  \loigiai{
    Dãy số có 20 số hạng. Tổng là $(20 + 1) \times 20 : 2 = 210$.
  }
\end{ex}

% Câu 23
\begin{ex}
  Một cửa hàng bán được $240\text{ kg}$ gạo trong 3 ngày. Ngày thứ nhất bán được $70\text{ kg}$, ngày thứ hai bán nhiều hơn ngày thứ nhất $15\text{ kg}$. Hỏi ngày thứ ba cửa hàng bán được bao nhiêu ki-lô-mét gạo?
  \shortans{85}
  \loigiai{
    Ngày thứ hai bán được: $70 + 15 = 85\text{ kg}$. Ngày thứ ba bán được: $240 - 70 - 85 = 85\text{ kg}$.
  }
\end{ex}

% Câu 24
\begin{ex}
  Tìm một số tự nhiên, biết rằng nếu lấy số đó nhân với $8$ rồi trừ đi $45$ thì được kết quả là $155$.
  \shortans{25}
  \loigiai{
    Gọi số đó là $x$. Ta có: $x \times 8 - 45 = 155 \implies x \times 8 = 200 \implies x = 25$.
  }
\end{ex}

% Câu 25
\begin{ex}
  Tìm một số tự nhiên có hai chữ số, biết rằng khi viết thêm chữ số $2$ vào bên trái số đó ta được số mới gấp $9$ lần số ban đầu.
  \shortans{25}
  \loigiai{
    Gọi số ban đầu là $\overline{ab}$. Viết thêm chữ số 2 vào bên trái được $\overline{2ab} = 200 + \overline{ab}$. Theo đề bài: $200 + \overline{ab} = 9 \times \overline{ab} \implies 8 \times \overline{ab} = 200 \implies \overline{ab} = 25$.
  }
\end{ex}
